%قبل از شروع مسئله ترجیحاً چیزی ننویسید. اما اگر خواستید بنویسید، در قالب لیستی (فقرات و شمارش) بنویسید.
%برای اضافه کردن نمره، بعدی از شروع‌مسئله، [x] را اضافه کنید و به جای x  نمره مورد نظر را قرار دهید. برای زیرمسئله‌ها هم به همین صورت است.
%\شروع‌مسئله[10]
\شروع‌مسئله [25]
%\متن‌سیاه{سوال عملی اول:}\\
%\متن‌سیاه{دراپ‌تست}\\

دانیال قصد دارد در یک مسابقه‌ی تلویزیونی به صورت تلفنی  شرکت کند، در این مسابقه تعداد زیادی عدد وجود دارد، و فرد پس از دیدن این اعداد، باید به سوالات گزارشگر درباره‌ی این اعداد پاسخ دهند. هر سوال گزارشگر به صورت یک عدد است و دانیال باید بگوید که آیا در بین اعداد سه عدد وجود دارد که جمع آن‌ها برابر عدد داده‌شده شود یا نه.
 
 به دانیال کمک کنید که برنده شود تا سهمی از جایزه‌اش را به شما بدهد.

\متن‌سیاه{ورودی:}\\
شامل یک لیست از اعداد است که در یک خط آمده‌اند. اعداد می‌توانند منفی باشند و اندازه‌ی آن‌ها کمتر از
$10^10$
می‌باشد.

سپس یک عدد در خط بعدی می‌آید که نمایانگر تعداد سوالات گزارشگر است.

سپس در هر خط یک عدد می‌آید که یک سوال گزارشگر می‌باشد.

\متن‌سیاه{خروجی:}\\
به ازای هر سوال گزارشگر اگر به صورت جمع سه عدد نوشته می‌شد، در یک خط متن YES وگرنه NO را چاپ کنید.

\متن‌سیاه{محدودیت‌ها:}\\
استفاده از داده ساختار dictionary پایتون مجاز نمی باشد. برای مقایسه رشته ها اجازه استفاده از عملگر == را ندارید. همچنین اجازه استفاده از داده ساختارهایی که در کلاس درس تدریس نشده اند را ندارید.

\متن‌سیاه{ورودی و خروجی نمونه:}\\

تست اول
\begin{latin}
	\begin{tabularx}{\textwidth}{|X | X|}
		\hline
		\textbf{stdin} &
		\textbf{stdout}\\
		\hline
		1 2 1000 2000
 & NO\\
		12 & NO \\
4 & NO\\
5 & NO \\
6 & NO \\
7 & NO \\
2 & YES \\
1 & NO \\
1003 & YES \\
3003 & YES \\
3002 & YES \\
3001 & YES \\
2003 & NO \\
3000 & \\
		\hline
	\end{tabularx}
\end{latin}


\solution{

ایده‌ی اصلی این سوال ذخیره‌ی جمع جفت عددها در یک جدول درهم‌سازی و سپس به ازای هر کوئری $q$ جست‌وجو کردن $q-arr[i]$ در این جدول است که $arr[i]$ برابر $i$مین عدد داده‌شده بین اعداد گزارشگر است. چالشی که در این زمینه‌ وجود دارد، اطمینان از متمایز بودن اعداد است.

برای حل این مشکل کافیست به ازای هر جفت عدد یک آرایه به جمع آن‌ها در جدول درهم سازی نسبت دهیم و در آن همه‌ی جفت اندیس‌های با آن جمع را نگه داریم. همه‌ی جمع جفت عدد‌ها و اندیسشان را در یک hashset ذخیره ‌کنیم. (برای هش کردن فقط از جمع استفاده می‌کنیم) هر عملیات ذخیره‌سازی از $O(1)$ و تعداد این عملیات‌ها از $O(n^2)$ می‌باشد.

برای هر کوئری‌ها نیز با حالت بندی کردن عدد سوم $c$ مقدار $q-c$ را در این جدول جست‌وجو کرده و با مقایسه‌ی اندیس‌های ذخیره‌شده از متمایز بودن سه عدد مطمئن می‌شویم. به راحتی می‌توانید مشاهده کنید که اگر تعداد اندیس‌های ذخیره شده به ازای یک جمع بیشتر از یک باشد، آنگاه حتما سه عدد متمایز پیدا می‌شود. پس به ازای هر کوئری می‌توان در $O(n)$ معتبر بودن آن را بررسی کرد. 

کد نمونه (با تشکر از کیان باختری):
\\

\definecolor{codegreen}{rgb}{0,0.6,0}
\definecolor{codegray}{rgb}{0.5,0.5,0.5}
\definecolor{codepurple}{rgb}{0.58,0,0.82}
\definecolor{backcolour}{rgb}{0.95,0.95,0.92}

\lstdefinestyle{mystyle}{
    backgroundcolor=\color{backcolour},   
    commentstyle=\color{codegreen},
    keywordstyle=\color{magenta},
    numberstyle=\tiny\color{codegray},
    stringstyle=\color{codepurple},
    basicstyle=\ttfamily\footnotesize,
    breakatwhitespace=false,         
    breaklines=true,                 
    captionpos=b,                    
    keepspaces=true,                 
    numbers=left,                    
    numbersep=5pt,                  
    showspaces=false,                
    showstringspaces=false,
    showtabs=false,                  
    tabsize=2
}

\lstset{style=mystyle}

\lr{\lstinputlisting[language=python]{src/kian.py}}



}

%}
% صورت اصلی مسئله را در این‌جا بنویسید.
%اگر مسئله چند قسمت داشت به این صورت عمل کنید:
%\شروع{زیرمسائل}
%\زیرمسئله[0] % مسئلۀ ۱آ
%\زیرمسئله % مسئلۀ ۱ب
%\زیرمسئله % مسئلۀ ۱پ
%\زیرمسئله % مسئلۀ ۱ت
%\پایان{زیرمسائل}



%%%%%%%%%%%%%%%%%%%%%%%%%%%%%%%%%%%%%%%%%%%%%%%%%%%%%
% برخی دستورات مفید برای تمرین‌ها. %
%%%%%%%%%%%%%%%%%%%%%%%%%%%%%%%%%%%%%%%%%%%%%%%%%%%%%

% فرمول‌های نمونه:
%$f(x) = \Theta(x)$
%$T(x, y) \leq \log(x) + 2^y + \binom{2n}{n}$
% {\tt code\_function}


% لیست موارد بدون شماره‌گذاری:
%\شروع{فقرات}
%    \فقره مورد اول 
%        \شروع{فقرات}
%            \فقره مورد اول داخلی 
%                \شروع{فقرات}
%                    \فقره اولین مورد خیلی داخلی 
%                    \فقره دومین مورد خیلی داخلی 
%                \پایان{فقرات}
%            \فقره مورد دوم داخلی 
%        \پایان{فقرات}
%    \فقره مورد دوم 
%\پایان{فقرات}

% لیست موارد شماره‌گذاری شده:
%\شروع{شمارش}
%    \فقره یک مورد 
%    \فقره یک مورد دیگر 
%    \فقره باز هم یک مورد دیگر
%\پایان{شمارش}

% برای نوشتن یک فرمول چند خطی که یک شماره بخوره. برای حالت‌های دیگر به سایت زیر مراجعه کنید:
% https://www.overleaf.com/learn/latex/Aligning_equations_with_amsmath
%\begin{align} 
%    \begin{split}
%        (x+y)^3 &= (x+y)^2(x+y) \\
%                &= (x^2+2xy+y^2)(x+y) \\
%                &= (x^3+2x^2y+xy^2) + (x^2y+2xy^2+y^3) \\
%                &= x^3+3x^2y+3xy^2+y^3
%    \end{split}                                 
%\end{align}

% نوشتن ماتریس و آرایه دوبعدی
%\begin{align}
%    A = 
%    \begin{bmatrix}
%        A_{11} & A_{21} \\
%        A_{21} & A_{22}
%    \end{bmatrix}
%\end{align}

% اضافه کردن تصویر با شرح. برای اضافه کردن تصویر بدون شرح، فقط دستور درج‌تصویر را استفاده کنید.
%\شروع{شکل}[h!]
%\تنظیم‌ازوسط
%  \درج‌تصویر[پهنا=0.5\پهنای‌متن]{img.jpg}
%\شرح{شرح تصویر}
%\برچسب{تصویر:۱}
%\پایان{شکل}
%\پایان{زیرمسائل}

%کد زیر را ببینید.
%%\lstinputlisting{source_filename.py}
%\begin{latin}
%\begin{lstlisting}
%# Import the modules
%import sys
%import random
%
%ans = True
%
%while ans:
%question = raw_input("Ask the magic 8 ball a question: (press enter to quit) ")
%
%answers = random.randint(1,8)
%
%if question == "":
%sys.exit()
%
%elif answers == 1:
%print "It is certain"
%
%elif answers == 2:
%print "Outlook good"
%
%elif answers == 3:
%print "You may rely on it"
%
%elif answers == 4:
%print "Ask again later"
%
%elif answers == 5:
%print "Concentrate and ask again"
%
%elif answers == 6:
%print "Reply hazy, try again"
%
%elif answers == 7:
%print "My reply is no"
%
%elif answers == 8:
%print "My sources say no"
%\end{lstlisting}
%\end{latin}
